\documentclass[12pt]{article}
\usepackage{sbc-template}
\usepackage[T1]{fontenc}
\usepackage[utf8]{inputenc}
\usepackage[brazil]{babel}
\usepackage{url}
\usepackage{booktabs}
\usepackage{array}
\titlespacing*{\section}{0pt}{12pt}{6pt}
\titlespacing*{\subsection}{0pt}{8pt}{4pt}
\sloppy
\title{DebugMesh: Distribui\c{c}\~ao de An\'alise de Logs com RPC, Rel\'ogios L\'ogicos e Elei\c{c}\~ao de L\'ider}
\author{Gustavo Castro, Vitor De Meira Gomes, Mateus R Fernandes,\\{}Felipe Barros, Arthur Carvalho}
\address{Ci\^encia da Computa\c{c}\~ao -- PUC Minas, Cora\c{c}\~ao Eucar\'istico\\Disciplina: Computa\c{c}\~ao Distribu\'ida -- Prof. Luigi D. C. Soares\\2026/2}
\begin{document}
\maketitle
\begin{resumo}
Este artigo apresenta a primeira entrega do DebugMesh, um prot\'otipo de servi\c{c}o distribu\'ido para an\'alise de logs de pipelines de integra\c{c}\~ao cont\'inua. O sistema divide registros em lotes, distribui o trabalho entre processos independentes por gRPC, aplica regras de diagn\'ostico e consolida evid\^encias. A implementa\c{c}\~ao inclui rel\'ogios l\'ogicos de Lamport, elei\c{c}\~ao simplificada de l\'ider e tratamento expl\'icito de indisponibilidade. Descrevem-se arquitetura, mecanismos, resultados dos testes e limita\c{c}\~oes.
\end{resumo}

\section{Introdu\c{c}\~ao e problema}
Falhas em pipelines de integra\c{c}\~ao cont\'inua podem envolver compila\c{c}\~ao, testes e infraestrutura, gerando extensos registros para investiga\c{c}\~ao. O DebugMesh foi proposto para demonstrar como lotes independentes de logs podem ser classificados em paralelo por processos conectados por rede. Sua entrada \'e um arquivo de logs; a sa\'ida \'e uma cole\c{c}\~ao de evid\^encias, categorias, hip\'oteses de causa, sugest\~oes de verifica\c{c}\~ao e identificadores dos workers participantes. O sistema n\~ao modifica c\'odigo automaticamente e n\~ao depende de uma API externa de intelig\^encia artificial.

Nesta entrega, o foco \'e o uso de tr\^es mecanismos de sistemas distribu\'idos: comunica\c{c}\~ao RPC, ordena\c{c}\~ao causal por rel\'ogios l\'ogicos e elei\c{c}\~ao de l\'ider. O trabalho foi desenvolvido em Python e usa gRPC com Protocol Buffers para estabelecer contratos entre os n\'os.

\section{Arquitetura e protocolo}
O cliente submete os registros ao coordenador ativo, que cria tarefas identificadas por \texttt{analysis\_id} e \texttt{task\_id}. Cada lote segue por \texttt{AssignTask} para um dos workers; eles processam as linhas e retornam resultados estruturados. O coordenador agrega os diagn\'osticos, preserva os resultados em mem\'oria e exp\~oe \texttt{GetAnalysisStatus} para consulta.

\begin{center}
\small
\begin{tabular}{c}
\fbox{Cliente: SubmitLogs / GetAnalysisStatus}\\[5pt]
$\downarrow$ gRPC\\[3pt]
\fbox{Coordenador ativo: fila + agregador}\\[5pt]
$\downarrow$ AssignTask, em paralelo\\[3pt]
\fbox{Worker A}\quad\fbox{Worker B}\quad\fbox{Worker C}\\[5pt]
$\uparrow$ resultados com evid\^encias e timestamps
\end{tabular}
\end{center}

O prot\'otipo tamb\'em disponibiliza um segundo processo eleg\'ivel \`a coordena\c{c}\~ao para demonstra\c{c}\~ao de elei\c{c}\~ao. Cada candidato conhece previamente o endere\c{c}o e a prioridade dos demais. O contrato em \texttt{debugmesh.proto} especifica \texttt{Envelope} com identificador da requisi\c{c}\~ao, emissor e timestamp l\'ogico. Os workers retornam \texttt{Finding}, incluindo categoria e evid\^encias. A execu\c{c}\~ao local utiliza dois coordenadores candidatos e tr\^es workers; para demonstrar distribui\c{c}\~ao de lotes, os tr\^es workers recebem subtarefas distintas.

\section{Mecanismos distribu\'idos}
\subsection{Comunica\c{c}\~ao RPC e paralelismo}
O cliente usa \texttt{SubmitLogs} para registrar a an\'alise e \texttt{GetAnalysisStatus} para consultar sua evolu\c{c}\~ao. Internamente, o coordenador despacha \texttt{AssignTask} de forma concorrente por gRPC. O tamanho dos lotes e os endere\c{c}os dos workers s\~ao configur\'aveis. Cada worker usa regras determin\'isticas para identificar falhas em compila\c{c}\~ao, testes, conex\~ao e exce\c{c}\~oes. O identificador do worker em cada resultado permite verificar a distribui\c{c}\~ao efetiva.

\subsection{Rel\'ogio l\'ogico de Lamport}
Cada n\'o possui contador protegido contra acessos concorrentes. Antes de emitir uma mensagem, incrementa seu rel\'ogio; ao receber uma mensagem com timestamp $T_m$, atualiza o contador para $L \leftarrow \max(L,T_m)+1$. Esses timestamps possibilitam registrar rela\c{c}\~oes causais entre eventos internos do DebugMesh, embora n\~ao reconstruam a ordem cronol\'ogica dos eventos que originaram os logs.

\subsection{Elei\c{c}\~ao de l\'ider}
Foi implementada uma variante did\'atica do algoritmo Bully. O candidato inicia uma elei\c{c}\~ao solicitando resposta de outros candidatos com prioridade superior. Se nenhum responder no prazo, declara-se l\'ider e anuncia sua escolha aos pares. O coordenador que sabe n\~ao ser l\'ider recusa submiss\~oes com \texttt{FAILED\_PRECONDITION}; processos secund\'arios monitoram a conectividade do l\'ider para iniciar nova elei\c{c}\~ao diante de sua queda. N\~ao existe qu\'orum nem consenso, portanto parti\c{c}\~oes de rede podem produzir lideran\c{c}as conflitantes.

\section{Testes, desafios e limita\c{c}\~oes}
Os testes locais foram executados com \texttt{pytest} e servidores gRPC em processos independentes. Foram aprovados treze testes automatizados (seis de unidade e sete de integra\c{c}\~ao): ordena\c{c}\~ao causal de Lamport, mensagens atrasadas, concorr\^encia com 1.200 eventos, classifica\c{c}\~ao, serializa\c{c}\~ao Protobuf, chamadas gRPC entre tr\^es workers, valida\c{c}\~ao de entradas e erros, indisponibilidade de worker, an\'alises concorrentes sem esgotamento do pool, elei\c{c}\~ao, sucess\~ao do l\'ider e rejei\c{c}\~ao de an\'uncios ou prioridades inconsistentes. A valida\c{c}\~ao automatizada no GitHub Actions (execu\c{c}\~ao 41158882709, em 30/10/2026) aprovou a su\'ite nas vers\~oes Python 3.11 e 3.12, al\'em de uma execu\c{c}\~ao em \mbox{Docker Compose}: as 12 linhas de amostra geraram seis diagn\'osticos distribu\'idos entre tr\^es workers. Ap\'os interromper o coordenador de maior prioridade, o secund\'ario assumiu a lideran\c{c}a e concluiu nova an\'alise. A demonstra\c{c}\~ao em cont\^eineres na rede virtual foi autorizada pelo professor; n\~ao se afirma ter executado o teste em computadores f\'isicos distintos.



O principal desafio foi definir contratos que permitissem execu\c{c}\~ao independente entre os servi\c{c}os sem confundir agentes de diagn\'ostico com coordena\c{c}\~ao. O estado de an\'alises \'e mantido em mem\'oria: caso o l\'ider caia, as an\'alises previamente aceitas n\~ao s\~ao recuperadas pelo sucessor. Reatribui\c{c}\~ao autom\'atica, replica\c{c}\~ao, autentica\c{c}\~ao e toler\^ancia a parti\c{c}\~oes ficam fora do escopo desta entrega. O ambiente de execu\c{c}\~ao utilizado foi o Docker Compose, conforme a modalidade aceita pelo professor.

\section{Participa\c{c}\~ao dos integrantes e considera\c{c}\~oes finais}
As responsabilidades da equipe foram organizadas em cinco frentes, com divis\~ao entre os integrantes:

\begin{itemize}
\setlength{\itemsep}{1pt}
\setlength{\parsep}{0pt}
\item \textbf{Gustavo Castro --- arquitetura e integra\c{c}\~ao:} defini\c{c}\~ao de requisitos, modelagem da arquitetura distribu\'ida, integra\c{c}\~ao dos componentes e consolida\c{c}\~ao da documenta\c{c}\~ao e do relat\'orio.
\item \textbf{Vitor De Meira Gomes --- comunica\c{c}\~ao gRPC:} contratos Protocol Buffers, comunica\c{c}\~ao cliente--coordenador, servi\c{c}os RPC, interface de linha de comando e valida\c{c}\~ao de requisi\c{c}\~oes.
\item \textbf{Mateus R Fernandes --- workers e paralelismo:} implementa\c{c}\~ao dos workers, divis\~ao dos logs em lotes, distribui\c{c}\~ao concorrente, regras de diagn\'ostico, agrega\c{c}\~ao de resultados e testes de processamento.
\item \textbf{Felipe Barros --- rel\'ogios l\'ogicos:} mecanismo de Lamport, timestamps nas mensagens, instrumenta\c{c}\~ao dos eventos de envio e recebimento e testes de causalidade e concorr\^encia.
\item \textbf{Arthur Carvalho --- elei\c{c}\~ao e valida\c{c}\~ao:} algoritmo Bully, detec\c{c}\~ao de falhas e sucess\~ao de lideran\c{c}a, Docker Compose, testes de failover e automa\c{c}\~ao pelo GitHub Actions.
\end{itemize}

A integra\c{c}\~ao, os testes e a revis\~ao final foram considerados responsabilidades compartilhadas. O c\'odigo-fonte, as instru\c{c}\~oes de execu\c{c}\~ao e as evid\^encias est\~ao dispon\'iveis no reposit\'orio \url{https://github.com/GhrCastro/TP01-Computacao-Distribuida}, na branch \texttt{release/primeira-entrega}. A execu\c{c}\~ao automatizada verificada \'e a de n\'umero \texttt{41158882709}.

O prot\'otipo demonstra comunica\c{c}\~ao, distribui\c{c}\~ao de trabalho e coordena\c{c}\~ao entre processos. Como continuidade, prop\~oem-se persist\^encia, replica\c{c}\~ao, retries com idempot\^encia e experimentos de desempenho.

\begin{thebibliography}{1}
\bibitem{lamport} L. Lamport. Time, Clocks, and the Ordering of Events in a Distributed System. \textit{Communications of the ACM}, 21(7), 1978.
\bibitem{grpc} gRPC Authors. gRPC Documentation. \url{https://grpc.io/docs/}.
\end{thebibliography}
\end{document}
